% preamble.tex -- shared preamble of the comparison note (orchestrator, 31 October 2026).
% Parts must not define macros; everything not listed here is written out explicitly.
\documentclass[11pt,a4paper]{article}
\usepackage[margin=2.9cm]{geometry}
\usepackage{amsmath,amssymb,amsthm,mathrsfs}
\usepackage[T1]{fontenc}
\usepackage{microtype}
\usepackage{booktabs}
\usepackage[colorlinks=true,linkcolor=blue,citecolor=blue]{hyperref}

% --- macros of part_i/main.tex, same names and meanings ---
\newcommand{\R}{\mathbb{R}}
\newcommand{\Z}{\mathbb{Z}}
\newcommand{\E}{\mathbb{E}}
\newcommand{\A}{\mathfrak{A}}
\newcommand{\G}{\mathcal{G}}
\newcommand{\CE}{\mathcal{C}_{\mathrm E}}
\newcommand{\CH}{\mathcal{C}_{\mathrm H}}
\newcommand{\HH}{\mathcal{H}}
\newcommand{\Om}{\Omega}
\newcommand{\ip}[2]{\langle #1,#2\rangle}
\DeclareMathOperator{\supp}{supp}
\DeclareMathOperator{\dist}{dist}
\DeclareMathOperator{\spec}{spec}
% --- macros of this note ---
\newcommand{\UGJS}{U_{\rm GJS}}
\newcommand{\GJS}[1]{\lVert #1\rVert_{\mathrm{GJS}}}   % the printed norm (1.1.7)-(1.1.9), eq. (1) of the record
\newcommand{\ev}[2]{\langle #1\rangle_{#2}}            % \ev{F}{h} = \int F dq_h
\DeclareMathOperator{\Cov}{Cov}                        % \Cov_h(F,G) = \ev{FG}{h}-\ev{F}{h}\ev{G}{h}
\newcommand{\Uc}{\mathcal U}                           % \Uc_h = \lambda\int h :P(\Phi): dz, the record's A3
\newcommand{\dmu}{d\mu_0}                              % free Euclidean measure
\newcommand{\dphi}{d\varphi_0}                         % free time-zero measure

% --- numbered statements ---
\theoremstyle{plain}
\newtheorem{dtheorem}{Theorem}
\renewcommand{\thedtheorem}{D.\arabic{dtheorem}}
\newtheorem{dlemma}[dtheorem]{Lemma}
\newtheorem{auxlemma}{Lemma}
\renewcommand{\theauxlemma}{N.\arabic{auxlemma}}
\newtheorem{auxprop}[auxlemma]{Proposition}
\newtheorem{auxcor}[auxlemma]{Corollary}
\theoremstyle{definition}
\newtheorem{auxdef}[auxlemma]{Definition}
\newtheorem{auxrem}[auxlemma]{Remark}
\newtheorem*{gapnote}{GAP}
